\documentclass[10pt,oneside,openany]{memoir}

\iffalse
%font style 1
\usepackage[p,osf]{scholax}
\usepackage{amsmath,amsthm,amssymb}
\usepackage[scaled=1.075,ncf,vvarbb]{newtxmath}
\fi

%font style 2
\usepackage[T1]{fontenc}
\usepackage{amsmath}
\usepackage[fixamsmath]{mathtools}  % Extension to amsmath
\usepackage{amsthm}
\usepackage{amssymb}
\renewcommand*{\mathbf}[1]{\varmathbb{#1}}
\usepackage{newpxtext}
\usepackage{eulerpx}
\usepackage{mathrsfs}

\usepackage{eucal}
\usepackage{datetime}
    \newdateformat{specialdate}{\THEYEAR\ \monthname\ \THEDAY}
\usepackage[margin=1.5in]{geometry}
\usepackage{fancyhdr}
    \pagestyle{fancy}
    \fancyhf{}
    \cfoot{\scriptsize\thepage}
    \fancyhead[R]{\scalebox{0.8}{\nouppercase{\rightmark}}}
    \fancyhead[L]{\scalebox{0.8}{\nouppercase{\leftmark}}}

    \fancypagestyle{plain}{%
        \fancyhf{}
        \cfoot{\scriptsize\thepage}
        \fancyhead[R]{\scalebox{0.8}{\phantom{a}}}
        \fancyhead[L]{\scalebox{0.8}{Applied Analysis}}
    }
\usepackage{keytheorems}
    \newkeytheoremstyle{def}{
    inherit-style=definition,
    spaceabove=1.5em,
    spacebelow=1.5em,
    postheadspace=0.5em,
    }
    \newkeytheoremstyle{thm}{
    spaceabove=1.5em,
    spacebelow=1.5em,
    postheadspace=0.5em,
    }
    \newkeytheorem{theorem}[
    name = Theorem,
    style=thm,
    numberwithin = section,
    ]
    \newkeytheorem{theorem*}[
    name=Theorem,
    style=thm,
    numbered=false
    ]
    \newkeytheorem{proposition}[
    name = Proposition,
    style=thm,
    numberwithin = section,
    sibling=theorem
    ]
    \newkeytheorem{lemma}[
    name = Lemma,
    style=thm,
    numberwithin = section,
    sibling=theorem
    ]
    \newkeytheorem{corollary}[
    name = Corollary,
    style=thm,
    numberwithin = section,
    sibling=theorem
    ]
    \newkeytheorem{definition}[
    name = Definition,
    style=def,
    numberwithin = section,
    style = definition
    ]
    \newkeytheorem{example}[
    name = Example,
    style=def,
    numberwithin = section,
    style = definition
    ]
    \newkeytheorem{remark}[
    name=Remark,
    style=remark, 
    numbered=false  
    ]
    \newkeytheorem{solution}[
    name=Solution,
    style=remark, 
    numbered=false  
    ]
    \newkeytheorem{question}[
    name=Question,
    style=remark, 
    numbered=false  
    ]
    \newkeytheorem{exercise}[
    name = Exercise,
    style=thm
    ]
    \newkeytheorem{axiom}[name = Axiom]
\usepackage{enumitem}
\usepackage[explicit]{titlesec}
    \titleformat{\chapter}[display]
    {\bfseries\LARGE\raggedright}
    {Chapter {\thechapter}}
    {.1ex minus .1ex}
    {\Large #1}
    \titlespacing{\chapter}
    {3pc}{*3}{40pt}[3pc]

    \titleformat{\section}[block]
    {\normalfont\bfseries\large}
    {}{0pt}
    {\fbox{\strut \S\ \thesection.\ #1}}
    \titlespacing{\section}
    {0pt}{3ex plus .1ex minus .2ex}{3ex plus .1ex minus .2ex}
    \setlength{\fboxsep}{4pt}   % padding inside the box
    \setlength{\fboxrule}{0.6pt} % border thickness
    
    \titleformat{name=\section,numberless}[block]
    {\normalfont\bfseries\large\raggedright}
    {}{0pt}
    {\fbox{\strut #1}}
\usepackage[utf8x]{inputenc}
\usepackage{tikz}
\usepackage{tikz-cd}
\usepackage{float}
\usetikzlibrary{patterns}
\usetikzlibrary{positioning,arrows.meta}
\usetikzlibrary{decorations.pathreplacing}
\usepackage{wasysym}
\usepackage{hyperref}
\usepackage{ulem}
\usepackage{pgfplots}
\usetikzlibrary{patterns}
\usepgfplotslibrary{fillbetween}
\pgfplotsset{compat=1.18}
\usepackage{xcolor}
\usepackage{matlab-prettifier}
\lstdefinestyle{Matlab-console}{
  basicstyle=\ttfamily\small,
  numbers=none,
  frame=single,
  breaklines=true,
  showstringspaces=false 
}
\definecolor{matlabbg}{HTML}{F2F2F2}
\definecolor{matlabln}{HTML}{6E6E6E}

\lstdefinestyle{MatlabEditorGrey}{
  style=Matlab-editor,
  backgroundcolor=\color{matlabbg},
  numbers=left,
  numberstyle=\color{matlabln}\ttfamily\scriptsize,
  stepnumber=1,
  numbersep=10pt,
  xleftmargin=2.2em,     % space for line numbers
  frame=single,
  framerule=0pt,         % no border line (keep the box spacing)
  framesep=6pt,
  breaklines=true
}

\renewcommand{\headrulewidth}{0pt} % removes the top/header line
% (optional) if you also want no line above the footer:
\renewcommand{\footrulewidth}{0pt}
%%%%%%%%%%%%%%%%%%%%%%%%%%%%%%%%%%%%%%%%%%%%%%%%%%%%%%%%%%%%%
%%%%%%%%%%%%%%%%%%%%%%%%%%%%%%%%%%%%%%%%%%%%%%%%%%%%%%%%%%%%%
% ============================================================
% makros.tex — reorganized, full restored version
% ============================================================

% ============================================================
% Sha / Cyrillic support notes
% ============================================================

%to make the correct symbol for Sha
%\newcommand\cyr{%
%\renewcommand\rmdefault{wncyr}%
%\renewcommand\sfdefault{wncyss}%
%\renewcommand\encodingdefault{OT2}%
%\normalfont \selectfont} \DeclareTextFontCommand{\textcyr}{\cyr}

% ============================================================
% Math operators
% ============================================================

\DeclareMathOperator{\ab}{ab}
\DeclareMathOperator{\ad}{ad}
\DeclareMathOperator{\adj}{adj}
\DeclareMathOperator{\alg}{alg}
\DeclareMathOperator{\Alt}{Alt}
\DeclareMathOperator{\Ann}{Ann}
\DeclareMathOperator{\arith}{arith}
\DeclareMathOperator{\Aut}{Aut}
\DeclareMathOperator{\Be}{B}
\DeclareMathOperator{\Bd}{Bd}
\DeclareMathOperator{\card}{card}
\DeclareMathOperator{\Char}{char}
\DeclareMathOperator{\csp}{csp}
\DeclareMathOperator{\codim}{codim}
\DeclareMathOperator{\coker}{coker}
\DeclareMathOperator{\coh}{H}
\DeclareMathOperator{\compl}{compl}
\DeclareMathOperator{\conj}{conj}
\DeclareMathOperator{\cont}{cont}
\DeclareMathOperator{\Cov}{Cov}
\DeclareMathOperator{\crys}{crys}
\DeclareMathOperator{\Crys}{Crys}
\DeclareMathOperator{\cusp}{cusp}
\DeclareMathOperator{\diag}{diag}
\DeclareMathOperator{\diam}{diam}
\DeclareMathOperator{\Dom}{Dom}
\DeclareMathOperator{\disc}{disc}
\DeclareMathOperator{\dist}{dist}
\DeclareMathOperator{\Div}{div}
\DeclareMathOperator{\dR}{dR}
\DeclareMathOperator{\Eis}{Eis}
\DeclareMathOperator{\End}{End}
\DeclareMathOperator{\ev}{ev}
\DeclareMathOperator{\eval}{eval}
\DeclareMathOperator{\Eq}{Eq}
\DeclareMathOperator{\Ext}{Ext}
\DeclareMathOperator{\Fil}{Fil}
\DeclareMathOperator{\Fitt}{Fitt}
\DeclareMathOperator{\Frob}{Frob}
\DeclareMathOperator{\G}{G}
\DeclareMathOperator{\Gal}{Gal}
\DeclareMathOperator{\GL}{GL}
\DeclareMathOperator{\Gr}{Gr}
\DeclareMathOperator{\Graph}{Graph}
\DeclareMathOperator{\GSp}{GSp}
\DeclareMathOperator{\GUn}{GU}
\DeclareMathOperator{\Hom}{Hom}
\DeclareMathOperator{\id}{id}
\DeclareMathOperator{\Id}{Id}
\DeclareMathOperator{\Ik}{Ik}
\DeclareMathOperator{\IM}{Im}
\DeclareMathOperator{\Image}{im}
\DeclareMathOperator{\Ind}{Ind}
\DeclareMathOperator{\Inf}{inf}
\DeclareMathOperator{\ins}{ins}
\DeclareMathOperator{\inv}{inv}
\DeclareMathOperator{\Isom}{Isom}
\DeclareMathOperator{\J}{J}
\DeclareMathOperator{\Jac}{Jac}
\DeclareMathOperator{\lcm}{lcm}
\DeclareMathOperator{\length}{length}
\DeclareMathOperator*{\limit}{limit}
\DeclareMathOperator{\Log}{Log}
\DeclareMathOperator{\M}{M}
\DeclareMathOperator{\Mat}{Mat}
\DeclareMathOperator{\N}{N}
\DeclareMathOperator{\Nm}{Nm}
\DeclareMathOperator{\NIk}{N-Ik}
\DeclareMathOperator{\NSK}{N-SK}
\DeclareMathOperator{\new}{new}
\DeclareMathOperator{\obj}{obj}
\DeclareMathOperator{\old}{old}
\DeclareMathOperator{\ord}{ord}
\DeclareMathOperator{\Or}{O}
\DeclareMathOperator{\op}{op}
\DeclareMathOperator{\out}{out}
\DeclareMathOperator{\PGL}{PGL}
\DeclareMathOperator{\PGSp}{PGSp}
\DeclareMathOperator{\rank}{rank}
\DeclareMathOperator{\Ran}{Ran}
\DeclareMathOperator{\rel}{Rel}
\DeclareMathOperator{\Real}{Re}
\DeclareMathOperator{\RES}{res}
\DeclareMathOperator{\Res}{Res}
%\DeclareMathOperator{\Sha}{\textcyr{Sh}}
\DeclareMathOperator{\Sel}{Sel}
\DeclareMathOperator{\semi}{ss}
\DeclareMathOperator{\sgn}{sign}
\DeclareMathOperator{\SK}{SK}
\DeclareMathOperator{\SL}{SL}
\DeclareMathOperator{\SO}{SO}
\DeclareMathOperator{\Sp}{Sp}
\DeclareMathOperator{\Span}{span}
\DeclareMathOperator{\Spec}{Spec}
\DeclareMathOperator{\spin}{spin}
\DeclareMathOperator{\st}{st}
\DeclareMathOperator{\St}{St}
\DeclareMathOperator{\SUn}{SU}
\DeclareMathOperator{\supp}{supp}
\DeclareMathOperator{\Sup}{sup}
\DeclareMathOperator{\Sym}{Sym}
\DeclareMathOperator{\Tam}{Tam}
\DeclareMathOperator{\tors}{tors}
\DeclareMathOperator{\tr}{tr}
\DeclareMathOperator{\Tr}{Tr}
\DeclareMathOperator{\un}{un}
\DeclareMathOperator{\Un}{U}
\DeclareMathOperator{\val}{val}
\DeclareMathOperator{\vol}{vol}

% ============================================================
% General aliases and short macros
% ============================================================

\newcommand{\absgal}{\G_{\bbQ}}
\newcommand{\Loc}{L_{\operatorname{loc}}}
\renewcommand{\k}{\kappa}
\newcommand{\Ff}{F_{f}}
%\newcommand{\ts}{\,^{t}\!}
\newcommand{\lat}{\mathscr{L}}
\newcommand{\ov}[1]{\overline{#1}}
\newcommand{\up}{\upsilon}
\newcommand{\e}{\epsilon}
\newcommand{\tac}{\textasteriskcentered}

%mahesh macros
\newcommand{\tm}{\textrm}

\newcommand{\bone}{\mathbf{1}}
\newcommand{\sign}{\mathrm{sign}}
\newcommand{\eps}{\varepsilon}
\newcommand{\textui}[1]{\underline{\textit{#1}}}

%\newcommand{\ov}{\overline}
%\newcommand{\un}{\underline}
\newcommand{\fin}{\mathrm{fin}}
\newcommand{\nl}{\newline \mbox{}}
\newcommand{\h}[1]{\hspace{#1pt}}
\newcommand{\mtext}[1]{\hspace{6pt}\text{#1}\hspace{6pt}}
\newcommand{\lnorm}{\left\lVert}
\newcommand{\rnorm}{\right\rVert}
\newcommand{\ds}{\displaystyle}
\newcommand{\ts}{\textstyle}
\newcommand{\sfrac}[2]{{}^{#1}\mskip -5mu/\mskip -3mu_{#2}}

% ============================================================
% Category-style operators and contoured variants
% ============================================================

\DeclareMathOperator{\Sets}{S \mkern1.04mu e \mkern1.04mu t \mkern1.04mu s}
    \newcommand{\cSets}{\scalebox{1.02}{\contour{black}{$\Sets$}}}
    
\DeclareMathOperator{\Groups}{G \mkern1.04mu r \mkern1.04mu o \mkern1.04mu u \mkern1.04mu p \mkern1.04mu s}
    \newcommand{\cGroups}{\scalebox{1.02}{\contour{black}{$\Groups$}}}

\DeclareMathOperator{\TTop}{T \mkern1.04mu o \mkern1.04mu p}
    \newcommand{\cTop}{\scalebox{1.02}{\contour{black}{$\TTop$}}}

\DeclareMathOperator{\Htp}{H \mkern1.04mu t \mkern1.04mu p}
    \newcommand{\cHtp}{\scalebox{1.02}{\contour{black}{$\Htp$}}}

\DeclareMathOperator{\Mod}{M \mkern1.04mu o \mkern1.04mu d}
    \newcommand{\cMod}{\scalebox{1.02}{\contour{black}{$\Mod$}}}

\DeclareMathOperator{\Ab}{A \mkern1.04mu b}
    \newcommand{\cAb}{\scalebox{1.02}{\contour{black}{$\Ab$}}}

\DeclareMathOperator{\Rings}{R \mkern1.04mu i \mkern1.04mu n \mkern1.04mu g \mkern1.04mu s}
    \newcommand{\cRings}{\scalebox{1.02}{\contour{black}{$\Rings$}}}

\DeclareMathOperator{\ComRings}{C \mkern1.04mu o \mkern1.04mu m \mkern1.04mu R \mkern1.04mu i \mkern1.04mu n \mkern1.04mu g \mkern1.04mu s}
    \newcommand{\cComRings}{\scalebox{1.05}{\contour{black}{$\ComRings$}}}

\DeclareMathOperator{\hHom}{H \mkern1.04mu o \mkern1.04mu m}
    \newcommand{\cHom}{\scalebox{1.02}{\contour{black}{$\hHom$}}}

% ============================================================
% Mathcal
% ============================================================

\newcommand{\cA}{\mathcal{A}}
\newcommand{\cB}{\mathcal{B}}
\newcommand{\cC}{\mathcal{C}}
\newcommand{\cD}{\mathcal{D}}
\newcommand{\cE}{\mathcal{E}}
\newcommand{\cF}{\mathcal{F}}
\newcommand{\cG}{\mathcal{G}}
\newcommand{\cH}{\mathcal{H}}
\newcommand{\cI}{\mathcal{I}}
\newcommand{\cJ}{\mathcal{J}}
\newcommand{\cK}{\mathcal{K}}
\newcommand{\cL}{\mathcal{L}}
\newcommand{\cM}{\mathcal{M}}
\newcommand{\cN}{\mathcal{N}}
\newcommand{\cO}{\mathcal{O}}
\newcommand{\cP}{\mathcal{P}}
\newcommand{\cQ}{\mathcal{Q}}
\newcommand{\cR}{\mathcal{R}}
\newcommand{\cS}{\mathcal{S}}
\newcommand{\cT}{\mathcal{T}}
\newcommand{\cU}{\mathcal{U}}
\newcommand{\cV}{\mathcal{V}}
\newcommand{\cW}{\mathcal{W}}
\newcommand{\cX}{\mathcal{X}}
\newcommand{\cY}{\mathcal{Y}}
\newcommand{\cZ}{\mathcal{Z}}

% ============================================================
% Mathscrl% 
%============================================================

\newcommand{\scra}{\mathscr{a}}
\newcommand{\scrA}{\mathscr{A}}
\newcommand{\scrb}{\mathscr{b}}
\newcommand{\scrB}{\mathscr{B}}
\newcommand{\scrc}{\mathscr{c}}
\newcommand{\scrC}{\mathscr{C}}
\newcommand{\scrd}{\mathscr{d}}
\newcommand{\scrD}{\mathscr{D}}
\newcommand{\scre}{\mathscr{e}}
\newcommand{\scrE}{\mathscr{E}}
\newcommand{\scrf}{\mathscr{f}}
\newcommand{\scrF}{\mathscr{F}}
\newcommand{\scrg}{\mathscr{g}}
\newcommand{\scrG}{\mathscr{G}}
\newcommand{\scrh}{\mathscr{h}}
\newcommand{\scrH}{\mathscr{H}}
\newcommand{\scri}{\mathscr{i}}
\newcommand{\scrI}{\mathscr{I}}
\newcommand{\scrj}{\mathscr{j}}
\newcommand{\scrJ}{\mathscr{J}}
\newcommand{\scrk}{\mathscr{k}}
\newcommand{\scrK}{\mathscr{K}}
\newcommand{\scrl}{\mathscr{l}}
\newcommand{\scrL}{\mathscr{L}}
\newcommand{\scrm}{\mathscr{m}}
\newcommand{\scrM}{\mathscr{M}}
\newcommand{\scrn}{\mathscr{n}}
\newcommand{\scrN}{\mathscr{N}}
\newcommand{\scro}{\mathscr{o}}
\newcommand{\scrO}{\mathscr{O}}
\newcommand{\scrp}{\mathscr{p}}
\newcommand{\scrP}{\mathscr{P}}
\newcommand{\scrq}{\mathscr{q}}
\newcommand{\scrQ}{\mathscr{Q}}
\newcommand{\scrr}{\mathscr{r}}
\newcommand{\scrR}{\mathscr{R}}
\newcommand{\scrs}{\mathscr{s}}
\newcommand{\scrS}{\mathscr{S}}
\newcommand{\scrt}{\mathscr{t}}
\newcommand{\scrT}{\mathscr{T}}
\newcommand{\scru}{\mathscr{u}}
\newcommand{\scrU}{\mathscr{U}}
\newcommand{\scrv}{\mathscr{v}}
\newcommand{\scrV}{\mathscr{V}}
\newcommand{\scrw}{\mathscr{w}}
\newcommand{\scrW}{\mathscr{W}}
\newcommand{\scrx}{\mathscr{x}}
\newcommand{\scrX}{\mathscr{X}}
\newcommand{\scry}{\mathscr{y}}
\newcommand{\scrY}{\mathscr{Y}}
\newcommand{\scrz}{\mathscr{z}}
\newcommand{\scrZ}{\mathscr{Z}}

% ============================================================
% Mathfrak (missing \fi)
% ============================================================

\newcommand{\fa}{\mathfrak{a}}
\newcommand{\fA}{\mathfrak{A}}
\newcommand{\fb}{\mathfrak{b}}
\newcommand{\fB}{\mathfrak{B}}
\newcommand{\fc}{\mathfrak{c}}
\newcommand{\fC}{\mathfrak{C}}
\newcommand{\fd}{\mathfrak{d}}
\newcommand{\fD}{\mathfrak{D}}
\newcommand{\fe}{\mathfrak{e}}
\newcommand{\fE}{\mathfrak{E}}
\newcommand{\ff}{\mathfrak{f}}
\newcommand{\fF}{\mathfrak{F}}
\newcommand{\fg}{\mathfrak{g}}
\newcommand{\fG}{\mathfrak{G}}
\newcommand{\fh}{\mathfrak{h}}
\newcommand{\fH}{\mathfrak{H}}
\newcommand{\fI}{\mathfrak{I}}
\newcommand{\fj}{\mathfrak{j}}
\newcommand{\fJ}{\mathfrak{J}}
\newcommand{\fk}{\mathfrak{k}}
\newcommand{\fK}{\mathfrak{K}}
\newcommand{\fl}{\mathfrak{l}}
\newcommand{\fL}{\mathfrak{L}}
\newcommand{\fm}{\mathfrak{m}}
\newcommand{\fM}{\mathfrak{M}}
\newcommand{\fn}{\mathfrak{n}}
\newcommand{\fN}{\mathfrak{N}}
\newcommand{\fo}{\mathfrak{o}}
\newcommand{\fO}{\mathfrak{O}}
\newcommand{\fp}{\mathfrak{p}}
\newcommand{\fP}{\mathfrak{P}}
\newcommand{\fq}{\mathfrak{q}}
\newcommand{\fQ}{\mathfrak{Q}}
\newcommand{\fr}{\mathfrak{r}}
\newcommand{\fR}{\mathfrak{R}}
\newcommand{\fs}{\mathfrak{s}}
\newcommand{\fS}{\mathfrak{S}}
\newcommand{\ft}{\mathfrak{t}}
\newcommand{\fT}{\mathfrak{T}}
\newcommand{\fu}{\mathfrak{u}}
\newcommand{\fU}{\mathfrak{U}}
\newcommand{\fv}{\mathfrak{v}}
\newcommand{\fV}{\mathfrak{V}}
\newcommand{\fw}{\mathfrak{w}}
\newcommand{\fW}{\mathfrak{W}}
\newcommand{\fx}{\mathfrak{x}}
\newcommand{\fX}{\mathfrak{X}}
\newcommand{\fy}{\mathfrak{y}}
\newcommand{\fY}{\mathfrak{Y}}
\newcommand{\fz}{\mathfrak{z}}
\newcommand{\fZ}{\mathfrak{Z}}

% ============================================================
% Mathbf
% ============================================================

\newcommand{\bfA}{\mathbf{A}}
\newcommand{\bfB}{\mathbf{B}}
\newcommand{\bfC}{\mathbf{C}}
\newcommand{\bfD}{\mathbf{D}}
\newcommand{\bfE}{\mathbf{E}}
\newcommand{\bfF}{\mathbf{F}}
\newcommand{\bfG}{\mathbf{G}}
\newcommand{\bfH}{\mathbf{H}}
\newcommand{\bfI}{\mathbf{I}}
\newcommand{\bfJ}{\mathbf{J}}
\newcommand{\bfK}{\mathbf{K}}
\newcommand{\bfL}{\mathbf{L}}
\newcommand{\bfM}{\mathbf{M}}
\newcommand{\bfN}{\mathbf{N}}
\newcommand{\bfO}{\mathbf{O}}
\newcommand{\bfP}{\mathbf{P}}
\newcommand{\bfQ}{\mathbf{Q}}
\newcommand{\bfR}{\mathbf{R}}
\newcommand{\bfS}{\mathbf{S}}
\newcommand{\bfT}{\mathbf{T}}
\newcommand{\bfU}{\mathbf{U}}
\newcommand{\bfV}{\mathbf{V}}
\newcommand{\bfW}{\mathbf{W}}
\newcommand{\bfX}{\mathbf{X}}
\newcommand{\bfY}{\mathbf{Y}}
\newcommand{\bfZ}{\mathbf{Z}}

\newcommand{\bfa}{\mathbf{a}}
\newcommand{\bfb}{\mathbf{b}}
\newcommand{\bfc}{\mathbf{c}}
\newcommand{\bfd}{\mathbf{d}}
\newcommand{\bfe}{\mathbf{e}}
\newcommand{\bff}{\mathbf{f}}
\newcommand{\bfg}{\mathbf{g}}
\newcommand{\bfh}{\mathbf{h}}
\newcommand{\bfi}{\mathbf{i}}
\newcommand{\bfj}{\mathbf{j}}
\newcommand{\bfk}{\mathbf{k}}
\newcommand{\bfl}{\mathbf{l}}
\newcommand{\bfm}{\mathbf{m}}
\newcommand{\bfn}{\mathbf{n}}
\newcommand{\bfo}{\mathbf{o}}
\newcommand{\bfp}{\mathbf{p}}
\newcommand{\bfq}{\mathbf{q}}
\newcommand{\bfr}{\mathbf{r}}
\newcommand{\bfs}{\mathbf{s}}
\newcommand{\bft}{\mathbf{t}}
\newcommand{\bfu}{\mathbf{u}}
\newcommand{\bfv}{\mathbf{v}}
\newcommand{\bfw}{\mathbf{w}}
\newcommand{\bfx}{\mathbf{x}}
\newcommand{\bfy}{\mathbf{y}}
\newcommand{\bfz}{\mathbf{z}}

% ============================================================
% Blackboard bold
% ============================================================

\newcommand{\bbA}{\mathbb{A}}
\newcommand{\bbB}{\mathbb{B}}
\newcommand{\bbC}{\mathbb{C}}
\newcommand{\bbD}{\mathbb{D}}
\newcommand{\bbE}{\mathbb{E}}
\newcommand{\bbF}{\mathbb{F}}
\newcommand{\bbG}{\mathbb{G}}
\newcommand{\bbH}{\mathbb{H}}
\newcommand{\bbI}{\mathbb{I}}
\newcommand{\bbJ}{\mathbb{J}}
\newcommand{\bbK}{\mathbb{K}}
\newcommand{\bbL}{\mathbb{L}}
\newcommand{\bbM}{\mathbb{M}}
\newcommand{\bbN}{\mathbb{N}}
\newcommand{\bbO}{\mathbb{O}}
\newcommand{\bbP}{\mathbb{P}}
\newcommand{\bbQ}{\mathbb{Q}}
\newcommand{\bbR}{\mathbb{R}}
\newcommand{\bbS}{\mathbb{S}}
\newcommand{\bbT}{\mathbb{T}}
\newcommand{\bbU}{\mathbb{U}}
\newcommand{\bbV}{\mathbb{V}}
\newcommand{\bbW}{\mathbb{W}}
\newcommand{\bbX}{\mathbb{X}}
\newcommand{\bbY}{\mathbb{Y}}
\newcommand{\bbZ}{\mathbb{Z}}

\newcommand{\Bmu}{\mbox{$\raisebox{-0.59ex}
  {$l$}\hspace{-0.18em}\mu\hspace{-0.88em}\raisebox{-0.98ex}{\scalebox{2}
  {$\color{white}.$}}\hspace{-0.416em}\raisebox{+0.88ex}
  {$\color{white}.$}\hspace{0.46em}$}{}}  %need graphicx and xcolor. this produces blackboard bold mu 

% ============================================================
% Greek-letter aliases
% ============================================================

\newcommand{\jota}{\jmath}

% ============================================================
% Matrices
% ============================================================

\newcommand{\bmat}{\left( \begin{matrix}}
\newcommand{\emat}{\end{matrix} \right)}

\newcommand{\bbmat}{\left[ \begin{matrix}}
\newcommand{\ebmat}{\end{matrix} \right]}

\newcommand{\pmat}{\left( \begin{smallmatrix}}
\newcommand{\epmat}{\end{smallmatrix} \right)}

\newcommand{\mat}[4]{\begin{pmatrix}{#1}&{#2}\\{#3}&{#4}\end{pmatrix}}

% ============================================================
% Residues and averaged integrals
% ============================================================

\newcommand{\res}[1]{\underset{#1}{\RES}\,}

\makeatletter
\newcommand*{\mint}[1]{%
  % #1: overlay symbol
  \mint@l{#1}{}%
}
\newcommand*{\mint@l}[2]{%
  % #1: overlay symbol
  % #2: limits
  \@ifnextchar\limits{%
    \mint@l{#1}%
  }{%
    \@ifnextchar\nolimits{%
      \mint@l{#1}%
    }{%
      \@ifnextchar\displaylimits{%
        \mint@l{#1}%
      }{%
        \mint@s{#2}{#1}%
      }%
    }%
  }%
}
\newcommand*{\mint@s}[2]{%
  % #1: limits
  % #2: overlay symbol
  \@ifnextchar_{%
    \mint@sub{#1}{#2}%
  }{%
    \@ifnextchar^{%
      \mint@sup{#1}{#2}%
    }{%
      \mint@{#1}{#2}{}{}%
    }%
  }%
}
\def\mint@sub#1#2_#3{%
  \@ifnextchar^{%
    \mint@sub@sup{#1}{#2}{#3}%
  }{%
    \mint@{#1}{#2}{#3}{}%
  }%
}
\def\mint@sup#1#2^#3{%
  \@ifnextchar_{%
    \mint@sup@sub{#1}{#2}{#3}%
  }{%
    \mint@{#1}{#2}{}{#3}%
  }%
}
\def\mint@sub@sup#1#2#3^#4{%
  \mint@{#1}{#2}{#3}{#4}%
}
\def\mint@sup@sub#1#2#3_#4{%
  \mint@{#1}{#2}{#4}{#3}%
}
\newcommand*{\mint@}[4]{%
  % #1: \limits, \nolimits, \displaylimits
  % #2: overlay symbol: -, =, ...
  % #3: subscript
  % #4: superscript
  \mathop{}%
  \mkern-\thinmuskip
  \mathchoice{%
    \mint@@{#1}{#2}{#3}{#4}%
        \displaystyle\textstyle\scriptstyle
  }{%
    \mint@@{#1}{#2}{#3}{#4}%
        \textstyle\scriptstyle\scriptstyle
  }{%
    \mint@@{#1}{#2}{#3}{#4}%
        \scriptstyle\scriptscriptstyle\scriptscriptstyle
  }{%
    \mint@@{#1}{#2}{#3}{#4}%
        \scriptscriptstyle\scriptscriptstyle\scriptscriptstyle
  }%
  \mkern-\thinmuskip
  \int#1%
  \ifx\\#3\\\else_{#3}\fi
  \ifx\\#4\\\else^{#4}\fi  
}
\newcommand*{\mint@@}[7]{%
  % #1: limits
  % #2: overlay symbol
  % #3: subscript
  % #4: superscript
  % #5: math style
  % #6: math style for overlay symbol
  % #7: math style for subscript/superscript
  \begingroup
    \sbox0{$#5\int\m@th$}%
    \sbox2{$#5\int_{}\m@th$}%
    \dimen2=\wd0 %
    % => \dimen2 = width of \int
    \let\mint@limits=#1\relax
    \ifx\mint@limits\relax
      \sbox4{$#5\int_{\kern1sp}^{\kern1sp}\m@th$}%
      \ifdim\wd4>\wd2 %
        \let\mint@limits=\nolimits
      \else
        \let\mint@limits=\limits
      \fi
    \fi
    \ifx\mint@limits\displaylimits
      \ifx#5\displaystyle
        \let\mint@limits=\limits
      \fi
    \fi
    \ifx\mint@limits\limits
      \sbox0{$#7#3\m@th$}%
      \sbox2{$#7#4\m@th$}%
      \ifdim\wd0>\dimen2 %
        \dimen2=\wd0 %
      \fi
      \ifdim\wd2>\dimen2 %
        \dimen2=\wd2 %
      \fi
    \fi
    \rlap{%
      $#5%
        \vcenter{%
          \hbox to\dimen2{%
            \hss
            $#6{#2}\m@th$%
            \hss
          }%
        }%
      $%
    }%
  \endgroup
}
\newcommand{\avg}{\mint{-}}
\makeatother

% ============================================================
% Comments and drafting helpers
% ============================================================

%Comments
\newcommand{\com}[1]{\vspace{5 mm}\par \noindent
\marginpar{\textsc{Comment}} \framebox{\begin{minipage}[c]{0.95
\textwidth} \tt #1 \end{minipage}}\vspace{5 mm}\par}

% ============================================================
% Arrows
% ============================================================

\newcommand{\hooktwoheadrightarrow}{%
  \hookrightarrow\mathrel{\mspace{-15mu}}\rightarrow
}

\makeatletter
\newcommand{\xhooktwoheadrightarrow}[2][]{%
  \lhook\joinrel
  \ext@arrow 0359\rightarrowfill@ {#1}{#2}%
  \mathrel{\mspace{-15mu}}\rightarrow
}
\makeatother

\makeatletter
\newcommand*{\da@rightarrow}{\mathchar"0\hexnumber@\symAMSa 4B }
\newcommand*{\da@leftarrow}{\mathchar"0\hexnumber@\symAMSa 4C }
\newcommand*{\xdashrightarrow}[2][]{%
  \mathrel{%
    \mathpalette{\da@xarrow{#1}{#2}{}\da@rightarrow{\,}{}}{}%
  }%
}
\newcommand{\xdashleftarrow}[2][]{%
  \mathrel{%
    \mathpalette{\da@xarrow{#1}{#2}\da@leftarrow{}{}{\,}}{}%
  }%
}
\newcommand*{\da@xarrow}[7]{%
  % #1: below
  % #2: above
  % #3: arrow left
  % #4: arrow right
  % #5: space left 
  % #6: space right
  % #7: math style 
  \sbox0{$\ifx#7\scriptstyle\scriptscriptstyle\else\scriptstyle\fi#5#1#6\m@th$}%
  \sbox2{$\ifx#7\scriptstyle\scriptscriptstyle\else\scriptstyle\fi#5#2#6\m@th$}%
  \sbox4{$#7\dabar@\m@th$}%
  \dimen@=\wd0 %
  \ifdim\wd2 >\dimen@
    \dimen@=\wd2 %   
  \fi
  \count@=2 %
  \def\da@bars{\dabar@\dabar@}%
  \@whiledim\count@\wd4<\dimen@\do{%
    \advance\count@\@ne
    \expandafter\def\expandafter\da@bars\expandafter{%
      \da@bars
      \dabar@ 
    }%
  }%  
  \mathrel{#3}%
  \mathrel{%   
    \mathop{\da@bars}\limits
    \ifx\\#1\\%
    \else
      _{\copy0}%
    \fi
    \ifx\\#2\\%
    \else
      ^{\copy2}%
    \fi
  }%   
  \mathrel{#4}%
}
\makeatother

% ============================================================
% Relation and delimiter spacing
% ============================================================

\renewcommand{\geq}{\geqslant}
\renewcommand{\leq}{\leqslant}
\newcommand{\midd}{\hspace{4pt}\middle|\hspace{4pt}}

% ============================================================
% Left Right Updated Deliminators
% ============================================================

\makeatletter
\NewDocumentCommand\Left{O{}}{%
    \IfValueTF{#1}{\notblank{#1}{\mathopen\bBigg@{#1}}{\left}}{\left}}
\NewDocumentCommand\Right{O{}}{%
    \IfValueTF{#1}{\notblank{#1}{\mathopen\bBigg@{#1}}{\right}}{\right}}
\makeatother

% ============================================================
% Restrictions
% ============================================================

\newcommand{\littletaller}{\mathchoice{\vphantom{\big|}}{}{}{}}
\newcommand\restr[2]{{% we make the whole thing an ordinary symbol
  \left.\kern-\nulldelimiterspace % automatically resize the bar with \right
  #1 % the function
  \littletaller % pretend it's a little taller at normal size
  \right|_{#2} % this is the delimiter
  }}

% ============================================================
% Chapter title formatting
% ============================================================

\newcommand{\chnum}{\titleformat
{\chapter} % command
[display] % shape
{\centering} % format
{\Huge \color{black} \shadowbox{\thechapter}} % label
{-0.5em} % sep (space between the number and title)
{\LARGE \color{black} \underline} % before-code
}

\newcommand{\chunnum}{\titleformat
{\chapter} % command
[display] % shape
{} % format
{} % label
{0em} % sep
{ \begin{flushright} \begin{tabular}{r}  \Huge \color{black}
} % before-code
[
\end{tabular} \end{flushright} \normalsize
] % after-code
}



\begin{document}
\begin{center}
    {\Large Math 502 \\[0.1in]Homework \#4}\\[.175in]
    {Name:} {\underline{Gianluca Crescenzo\hspace*{2in}}}\\[0.15in]
    \end{center}
    \vspace{4pt}
%%%%%%%%%%%%%%%%%%%%%%%%%%%%%%%%%%%%%%%%%%%%%%%%%%
\begin{exercise}[manual-num={4.0}]
Show the following three facts, for $\vec{\boldsymbol{f}}$ a constant vector in $\bfR^n$, $\vec{\boldsymbol{u}}$ a vector of variables in $\bfR^n$, and $A$ an $n \times n$ symmetric matrix:
\begin{enumerate}[label = (\alph*),itemsep=1pt,topsep=3pt]
    \item $\nabla(\vec{\boldsymbol{u}}^{\top}\vec{\boldsymbol{f}})=\vec{\boldsymbol{f}},$
    \item $\nabla(\vec{\boldsymbol{f}}^{\top}\vec{\boldsymbol{u}})=\vec{\boldsymbol{f}},$
    \item $\nabla(\vec{\boldsymbol{u}}^{\top}A\vec{\boldsymbol{u}})=2A\vec{\boldsymbol{u}}.$
\end{enumerate}
\end{exercise}
\begin{solution}
\phantom{a}
\begin{enumerate}[label=(\alph*),itemsep=4pt,topsep=4pt]
    \item Define $g:\bfR^n\to\bfR$ by $g(\vec{\boldsymbol{u}})=\vec{\boldsymbol{u}}^{\top}\vec{\boldsymbol{f}}$. Then:
        \begin{equation*}
        \begin{split}
            g(\vec{\boldsymbol{u}})=\sum_{i=1}^n u_i f_i.
        \end{split}
        \end{equation*}
    Hence:
        \begin{equation*}
        \begin{split}
            \nabla g(\vec{\boldsymbol{u}})
            &=\biggl(\frac{\partial g}{\partial u_1},\dots,\frac{\partial g}{\partial u_n}\biggr) \\
            &=\biggl(\frac{\partial}{\partial u_1}\sum_{i=1}^n u_i f_i,\dots,\frac{\partial}{\partial u_n}\sum_{i=1}^n u_i f_i\biggr) \\
            &= (f_1,\dots,f_n) \\
            &=\vec{\boldsymbol{f}}.
        \end{split}
        \end{equation*}

    \item Note that $\varphi(\vec{\boldsymbol{u}},\vec{\boldsymbol{v}}) = \vec{\boldsymbol{u}}^{\top}\vec{\boldsymbol{v}}$ defines an inner product on $\bfR^n$. Hence $\vec{\boldsymbol{u}}^{\top}\vec{\boldsymbol{v}} = \vec{\boldsymbol{v}}^{\top}\vec{\boldsymbol{u}}$. Therefore by part (a):
        \begin{equation*}
        \begin{split}
            \nabla(\vec{\boldsymbol{f}}^{\top}\vec{\boldsymbol{u}})
            &=\nabla(\vec{\boldsymbol{u}}^{\top}\vec{\boldsymbol{f}}) \\
            &=\vec{\boldsymbol{f}}.
        \end{split}
        \end{equation*}

    \item Define $h:\bfR^n\to\bfR$ by $h(\vec{\boldsymbol{u}})=\vec{\boldsymbol{u}}^{\top}A\vec{\boldsymbol{u}}$. Writing $A=(a_{ij})_{i,j=1}^n$, we have:
        \begin{equation*}
        \begin{split}
            h(\vec{\boldsymbol{u}})
            =\sum_{i=1}^n\sum_{j=1}^n u_i\,a_{ij}\,u_j.
        \end{split}
        \end{equation*}
    Note that for each $k=1,...,n$:
        \begin{equation*}
        \begin{split}
            \frac{\partial h}{\partial u_k}
            &= \frac{\partial}{\partial u_k}\sum_{i=1}^n\sum_{j=1}^n u_i a_{ij} u_j \\
            &= \sum_{j=1}^n a_{kj}u_j \;+\; \sum_{i=1}^n u_i a_{ik}.
        \end{split}
        \end{equation*}
    Hence:
        \begin{equation*}
        \begin{split}
            \nabla h(\vec{\boldsymbol{u}})
            = (A+A^{\top})\vec{\boldsymbol{u}}.
        \end{split}
        \end{equation*}
    Since $A$ is symmetric, $A^{\top}=A$, which gives:
        \begin{equation*}
        \begin{split}
            \nabla(\vec{\boldsymbol{u}}^{\top}A\vec{\boldsymbol{u}})
            = 2A\vec{\boldsymbol{u}}.
        \end{split}
        \end{equation*}
\end{enumerate}
\end{solution}
%%%%%%%%%%%%%%%%%%%%%%%%%%%%%%%%%%%%%%%%%%%%%%%%%%
\begin{exercise}[manual-num={4.1}]
    By considering the derivative and the second derivative of $g(\alpha)$, show that the parameter $\alpha$ in the Conjugate Gradient algorithm minimizes 
        \begin{equation*}
        \begin{split}
            g(\alpha) = \phi \bigl( \vec{\boldsymbol{u}}_{k-1} + \alpha \vec{\boldsymbol{p}}_{k-1}\bigr),
        \end{split}
        \end{equation*}
    where $\phi(\vec{\boldsymbol{u}}) = \frac{1}{2}\vec{\boldsymbol{u}}^\top A \vec{\boldsymbol{u}} - \vec{\boldsymbol{u}}^\top \vec{\boldsymbol{f}}$.
\end{exercise}
    \begin{solution}
        Observe that:
            \begin{equation*}
            \begin{split}
                g(\alpha)
                &= \phi(\vec{\boldsymbol{u}}_{k-1}+\alpha\vec{\boldsymbol{p}}_{k-1}) \\
                &= \frac12(\vec{\boldsymbol{u}}_{k-1}+\alpha\vec{\boldsymbol{p}}_{k-1})^{\top}
                A(\vec{\boldsymbol{u}}_{k-1}+\alpha\vec{\boldsymbol{p}}_{k-1})
                -(\vec{\boldsymbol{u}}_{k-1}+\alpha\vec{\boldsymbol{p}}_{k-1})^{\top}\vec{\boldsymbol{f}} \\
                &= \frac12\Bigl(\vec{\boldsymbol{u}}_{k-1}^{\top}A\vec{\boldsymbol{u}}_{k-1}
                +\alpha\,\vec{\boldsymbol{u}}_{k-1}^{\top}A\vec{\boldsymbol{p}}_{k-1}
                +\alpha\,\vec{\boldsymbol{p}}_{k-1}^{\top}A\vec{\boldsymbol{u}}_{k-1}
                +\alpha^2\,\vec{\boldsymbol{p}}_{k-1}^{\top}A\vec{\boldsymbol{p}}_{k-1}\Bigr)
                -\vec{\boldsymbol{u}}_{k-1}^{\top}\vec{\boldsymbol{f}}
                -\alpha\,\vec{\boldsymbol{p}}_{k-1}^{\top}\vec{\boldsymbol{f}} \\
                &= \frac12\,\vec{\boldsymbol{u}}_{k-1}^{\top}A\vec{\boldsymbol{u}}_{k-1}
                + \alpha\,\vec{\boldsymbol{p}}_{k-1}^{\top}A\vec{\boldsymbol{u}}_{k-1}
                + \frac12\,\alpha^2\,\vec{\boldsymbol{p}}_{k-1}^{\top}A\vec{\boldsymbol{p}}_{k-1}
                -\vec{\boldsymbol{u}}_{k-1}^{\top}\vec{\boldsymbol{f}}
                -\alpha\,\vec{\boldsymbol{p}}_{k-1}^{\top}\vec{\boldsymbol{f}}.
            \end{split}
            \end{equation*}
        The derivative is then:
            \begin{equation*}
            \begin{split}
                g'(\alpha)
                &= \vec{\boldsymbol{p}}_{k-1}^{\top}A\vec{\boldsymbol{u}}_{k-1}
                + \alpha\,\vec{\boldsymbol{p}}_{k-1}^{\top}A\vec{\boldsymbol{p}}_{k-1}
                -\vec{\boldsymbol{p}}_{k-1}^{\top}\vec{\boldsymbol{f}} \\
                &= \vec{\boldsymbol{p}}_{k-1}^{\top}(A\vec{\boldsymbol{u}}_{k-1}-\vec{\boldsymbol{f}})
                + \alpha\,\vec{\boldsymbol{p}}_{k-1}^{\top}A\vec{\boldsymbol{p}}_{k-1}.
            \end{split}
            \end{equation*}
        Setting $g'(\alpha) = 0$ gives:
            \begin{equation*}
            \begin{split}
                \alpha
                &= \frac{\vec{\boldsymbol{p}}_{k-1}^{\top}(\vec{\boldsymbol{f}}-A\vec{\boldsymbol{u}}_{k-1})}
                {\vec{\boldsymbol{p}}_{k-1}^{\top}A\vec{\boldsymbol{p}}_{k-1}} \\
                &= \frac{\vec{\boldsymbol{p}}_{k-1}^{\top}\vec{\boldsymbol{r}}_{k-1}}{\vec{\boldsymbol{p}}_{k-1}^{\top}A\vec{\boldsymbol{p}}_{k-1}}.
            \end{split}
            \end{equation*}
        The second derivate of $g(\alpha)$ is:
            \begin{equation*}
            \begin{split}
                g''(\alpha) = \vec{\boldsymbol{p}}_{k-1}^{\top}A\vec{\boldsymbol{p}}_{k-1}.
            \end{split}
            \end{equation*}
        If $A$ is symmetric and positive-definite, then $p_{k-1} \neq 0$ implies $g''(\alpha) = \vec{\boldsymbol{p}}_{k-1}^{\top}A\vec{\boldsymbol{p}}_{k-1} > 0$. Hence the $\alpha$ given above uniquely minimizes $g(\alpha)$.
    \end{solution}
%%%%%%%%%%%%%%%%%%%%%%%%%%%%%%%%%%%%%%%%%%%%%%%%%%
\newpage
\begin{exercise}[manual-num={4.3}]
The file \texttt{conjugate\_gradient\_broken.m} is a modified version of an m-file that implements the conjugate gradient algorithm (without preconditioning) to solve \(A\tilde{u}=\vec{f}\). The modification is that two lines of code were removed. Download the file from our Canvas site, and repair the file so that it successfully implements the conjugate gradient algorithm (as stated in the algorithm at the end of this assignment).
Rename the file \texttt{conjugate\_gradient.m}, and test it by running the command
\begin{verbatim}
>> u = conjugate_gradient(A,F,1.0e-5);
\end{verbatim}
where \(A\) is some large, symmetric, positive-definite matrix (defined as a sparse matrix in MATLAB), and \(F\) a column vector of the correct length.
\end{exercise}
    \begin{solution}
        The two lines of code removed were an initialization of $r$ and updating $u$ during the for-loop. Below is the corrected code.
        \begin{lstlisting}[style=MatlabEditorGrey]
function u = conjugate_gradient(A,f,tol)
%
% Example:
% x = conjugate_gradient(A,f,tol)
%
MAXITS = length(f);
u = 0*f;
r = f- A*u   %r was never initialized
p = r;
for k = 1:MAXITS
    w = A*p;
    alpha = (r'*r)/(p'*w);
    u = u + alpha*p;  %u was never updated
    rnew = r - alpha*w;
    if( norm(rnew) < tol )
        fprintf('Converged! its= %7.0f, tol=%10.3e\n', [k tol]);
        return;
    end
    beta = (rnew'*rnew)/(r'*r);
    p = rnew + beta*p;
    r = rnew;
end
fprintf('Caution: CG went to max iterations without converging!\n');
fprintf('MAXITS = %7.0f, tol =%10.3e\n', [MAXITS tol]);

end
        \end{lstlisting}
Running the code with the large sparse matrix given from \texttt{poisson.m} gives:
\newpage
    \begin{lstlisting}[style=Matlab-console]
>> load("/Users/gcrescenzo/Documents/MATLAB/vechw4.mat")
>> load("/Users/gcrescenzo/Documents/MATLAB/sparsemat.mat")
>> conjugate_gradient_broken

r =

   1.0e+03 *

   -0.9128
   -0.4837
   -0.5072
   -0.5320
   -0.5579
   -0.5851
   -0.6137
   -0.6436
   -0.6750
   -0.7079
   -0.7424
   -0.7787
   -0.8166
   -0.8565
   -0.8982
   -0.9420
   -0.9880
   -1.0362
   -1.0867
   -2.3673
   -0.4611
    0.0014
    0.0015
    0.0016
    0.0017
    0.0017
    0.0018
    0.0019
    0.0020
    0.0021
    0.0022
    0.0023
    0.0024
    0.0026
    0.0027
    0.0028
    0.0029
    0.0031
    0.0032
   -1.2538
   -0.4722
    0.0015
    0.0015
    0.0016
    0.0017
    0.0018
    0.0019
    0.0020
    0.0021
    0.0022
    0.0023
    0.0024
    0.0025
    0.0026
    0.0027
    0.0029
    0.0030
    0.0032
    0.0033
   -1.2840
   -0.4836
    0.0015
    0.0016
    0.0017
    0.0017
    0.0018
    0.0019
    0.0020
    0.0021
    0.0022
    0.0023
    0.0024
    0.0026
    0.0027
    0.0028
    0.0029
    0.0031
    0.0032
    0.0034
   -1.3150
   -0.4953
    0.0015
    0.0016
    0.0017
    0.0018
    0.0019
    0.0020
    0.0021
    0.0022
    0.0023
    0.0024
    0.0025
    0.0026
    0.0027
    0.0029
    0.0030
    0.0032
    0.0033
    0.0035
   -1.3467
   -0.5072
    0.0016
    0.0017
    0.0017
    0.0018
    0.0019
    0.0020
    0.0021
    0.0022
    0.0023
    0.0024
    0.0026
    0.0027
    0.0028
    0.0029
    0.0031
    0.0032
    0.0034
    0.0036
   -1.3791
   -0.5194
    0.0016
    0.0017
    0.0018
    0.0019
    0.0020
    0.0021
    0.0022
    0.0023
    0.0024
    0.0025
    0.0026
    0.0027
    0.0029
    0.0030
    0.0032
    0.0033
    0.0035
    0.0036
   -1.4123
   -0.5319
    0.0017
    0.0017
    0.0018
    0.0019
    0.0020
    0.0021
    0.0022
    0.0023
    0.0024
    0.0026
    0.0027
    0.0028
    0.0029
    0.0031
    0.0032
    0.0034
    0.0036
    0.0037
   -1.4464
   -0.5448
    0.0017
    0.0018
    0.0019
    0.0020
    0.0021
    0.0022
    0.0023
    0.0024
    0.0025
    0.0026
    0.0027
    0.0029
    0.0030
    0.0032
    0.0033
    0.0035
    0.0036
    0.0038
   -1.4812
   -0.5579
    0.0017
    0.0018
    0.0019
    0.0020
    0.0021
    0.0022
    0.0023
    0.0024
    0.0026
    0.0027
    0.0028
    0.0029
    0.0031
    0.0032
    0.0034
    0.0036
    0.0037
    0.0039
   -1.5169
   -0.5713
    0.0018
    0.0019
    0.0020
    0.0021
    0.0022
    0.0023
    0.0024
    0.0025
    0.0026
    0.0027
    0.0029
    0.0030
    0.0032
    0.0033
    0.0035
    0.0036
    0.0038
    0.0040
   -1.5535
   -0.5851
    0.0018
    0.0019
    0.0020
    0.0021
    0.0022
    0.0023
    0.0024
    0.0026
    0.0027
    0.0028
    0.0029
    0.0031
    0.0032
    0.0034
    0.0036
    0.0037
    0.0039
    0.0041
   -1.5909
   -0.5992
    0.0019
    0.0020
    0.0021
    0.0022
    0.0023
    0.0024
    0.0025
    0.0026
    0.0027
    0.0029
    0.0030
    0.0032
    0.0033
    0.0035
    0.0036
    0.0038
    0.0040
    0.0042
   -1.6292
   -0.6136
    0.0019
    0.0020
    0.0021
    0.0022
    0.0023
    0.0024
    0.0026
    0.0027
    0.0028
    0.0029
    0.0031
    0.0032
    0.0034
    0.0036
    0.0037
    0.0039
    0.0041
    0.0043
   -1.6685
   -0.6284
    0.0020
    0.0021
    0.0022
    0.0023
    0.0024
    0.0025
    0.0026
    0.0027
    0.0029
    0.0030
    0.0032
    0.0033
    0.0035
    0.0036
    0.0038
    0.0040
    0.0042
    0.0044
   -1.7087
   -0.6436
    0.0020
    0.0021
    0.0022
    0.0023
    0.0024
    0.0026
    0.0027
    0.0028
    0.0029
    0.0031
    0.0032
    0.0034
    0.0036
    0.0037
    0.0039
    0.0041
    0.0043
    0.0045
   -1.7499
   -0.6591
    0.0021
    0.0022
    0.0023
    0.0024
    0.0025
    0.0026
    0.0027
    0.0029
    0.0030
    0.0032
    0.0033
    0.0035
    0.0036
    0.0038
    0.0040
    0.0042
    0.0044
    0.0046
   -1.7920
   -0.6750
    0.0021
    0.0022
    0.0023
    0.0024
    0.0026
    0.0027
    0.0028
    0.0029
    0.0031
    0.0032
    0.0034
    0.0036
    0.0037
    0.0039
    0.0041
    0.0043
    0.0045
    0.0047
   -1.8352
   -0.6912
    0.0022
    0.0023
    0.0024
    0.0025
    0.0026
    0.0027
    0.0029
    0.0030
    0.0032
    0.0033
    0.0035
    0.0036
    0.0038
    0.0040
    0.0042
    0.0044
    0.0046
    0.0049
   -1.8794
   -1.4704
   -0.7975
   -0.8364
   -0.8772
   -0.9200
   -0.9649
   -1.0119
   -1.0613
   -1.1130
   -1.1673
   -1.2242
   -1.2840
   -1.3466
   -1.4123
   -1.4811
   -1.5534
   -1.6291
   -1.7086
   -1.7919
   -3.8092

Converged! its=      68, tol= 1.000e-05
>> 
\end{lstlisting}
    \end{solution}
%%%%%%%%%%%%%%%%%%%%%%%%%%%%%%%%%%%%%%%%%%%%%%%%%%
\begin{exercise}[manual-num={5.1}]
Prove that the ODE:
    \begin{equation*}
    \begin{split}
        u'(t) = \frac{1}{t^2 = u(t)^2}, \h9 \text{for }t \geq 1
    \end{split}
    \end{equation*}
has a unique solution for all time from any initial value $u(1) = \eta$.
\end{exercise}
    \begin{solution}
        Define $f(u,t) = \frac{1}{t^2 = u^2}$. Then:
            \begin{equation*}
            \begin{split}
                \frac{\partial f}{\partial u} =-\frac{2u}{(t^2 + u^2)^2}.
            \end{split}
            \end{equation*}
        On the given interval $t \geq 1$, note that:
            \begin{equation*}
            \begin{split}
                \left| \frac{\partial f}{\partial u} \right|
                & = \frac{2|u|}{(t^2 + u^2)^2} \\
                & \leq \frac{2|u|}{(1 + u^2)^2} \\
            \end{split}
            \end{equation*}
        The function $\frac{2|u|}{(1+u^2)^2}$ has a horizontal asymptote at 0 (because the degree of the denominator is larger than the degree of the numerator) and has no vertical asymptotes (because the denominator is never equal to zero). It must be the case that $\frac{2|u|}{(1+u^2)^2}$ is bounded. Hence $\frac{\partial f}{\partial u}$ is bounded, establishing that $f(u,t)$ is Lipschitz continuous on $t \geq 1$. Since $f(u,t)$ is Lipschitz continuous on $t \geq 1$, the initial value problem described above admits a unique solution.
    \end{solution}
%%%%%%%%%%%%%%%%%%%%%%%%%%%%%%%%%%%%%%%%%%%%%%%%%%
\begin{exercise}[manual-num={5.2}]
Let $f(u) = \log(u)$.
\begin{enumerate}[label = (\alph*),itemsep=1pt,topsep=3pt]
    \item Determine the best possible Lipschitz constant for this function over $2 \leq u < \infty$.
    \item Is $f(u)$ Lipschitz continuous over $0 < u < \infty$?
    \item Consider the initial value problem
        \begin{equation*}
        \begin{split}
            &u'(t) = \log(u(t)), \\
            &u(0) = 2.
        \end{split}
        \end{equation*}
    Explain why we know that this problem has a unique solution for all $t \geq 0$ based on the existence and uniqueness theory described in Section 5.2.1.
\end{enumerate}
\end{exercise}
    \begin{solution}
        (a) An immediate application of the Mean Value Theorem will show that $f(u) = \log(u)$ is Lipschitz. Let $x,y \in [2,\infty)$. Suppose we can find $L > 0$ such that $|f(x) - f(y)| \leq L|x-y|$. Then $\frac{|f(x) - f(y)|}{|x-y|} \leq L$. Taking the limit as $x \rightarrow y$ on both sides gives the inequality $|f'(y)| \leq L$; i.e., $|\frac{1}{y}| \leq L$. Taking the supremum over all $y \in [2,\infty)$ gives $\frac{1}{2} \leq L$. Thus $L = \frac{1}{2}$ is the best Lipschitz constant over the domain $[2,\infty)$.

        (b) A function is Lipschitz if and only if its derivative is bounded. Since $f'(u) = \frac{1}{x}$ is unbounded on $(0,\infty)$, $f(u) = \log(u)$ is not Lipschitz on $(0,\infty)$.

        (c) Note that $u(t) \geq 2$ for all $t \geq 0$ based off of the initial value given. Since $\log(x)$ is Lipschitz on the domain $[2,\infty)$, this initial value problem admits a unique solution for all $t \geq 0$.
    \end{solution}







\end{document}